\section{Filtered output norms and strong-converse exponents}\label{sec:renyi}
This section strengthens the order-one conclusions for the original Davies
channel. It does not assert the corresponding filtered-norm or exponent
identities for its complementary channel. The triangular Schatten inequality
below is Fang's theorem, not a result proved here. The coefficient triangle
is the one used in Section~2. The additional ingredient at its third vertex
is a standard consequence of Araki--Lieb--Thirring.

For a finite-dimensional completely positive map $\Psi$, not necessarily
trace preserving, define
\[
 \nu_p(\Psi)=\max_{\rho\ge0,\,\Tr\rho=1}\|\Psi(\rho)\|_p,
 \qquad 1\le p\le\infty.
\]
The maximum exists by compactness. For $H\ge0$ write
$\Phi_H(X)=H\Phi(X)H$.

\begin{lemma}[External triangular norm compression]\label{lem:fangnorm}
For every $r\ge2$ and compatible rectangular blocks,
\[
 \left\|\begin{pmatrix}U&0\\V&W\end{pmatrix}\right\|_r
 \le
 \left\|\begin{pmatrix}\|U\|_r&0\\\|V\|_r&\|W\|_r\end{pmatrix}\right\|_r.
\]
\end{lemma}
This is \cite[Theorem~3]{fang}, including its rectangular and operator-norm
endpoints. Its trace-power convexity argument is an external dependency of
the present extension.

\begin{lemma}[Controlled block factors]\label{lem:altnorm}
Let $C\ge0$ be an $n\times n$ scalar matrix and let $X_i$ be $d\times r$
matrices with $\|X_i\|_{2p}=1$, where $1\le p<\infty$. Then
\[
 \left\|\sum_{ij}C_{ij}|i\rangle\langle j|\otimes X_iX_j^*\right\|_p
 \le\|C\|_p.
\]
The analogous statement holds for $p=\infty$ with $\|X_i\|_\infty=1$.
\end{lemma}
\begin{proof}
Set $D=\operatorname{diag}(X_1,\ldots,X_n)$,
$A=C\otimes I_r$ and $B=D^*D$. The matrix on the left is $DAD^*$,
whose nonzero eigenvalues are those of $B^{1/2}AB^{1/2}$. The standard
Araki--Lieb--Thirring inequality \cite[Theorem~1]{alt} gives
\begin{align*}
 \Tr\bigl[(DAD^*)^p\bigr]
 &\le\Tr B^{p/2}A^p B^{p/2}\\
 &=\sum_i(C^p)_{ii}\Tr\bigl[(X_i^*X_i)^p\bigr]=\Tr C^p.
\end{align*}
Rectangular factors add zero eigenvalues and do not change this argument.
For $p=\infty$, $D$ is a contraction, so $\|DAD^*\|\le\|A\|=\|C\|$.
\end{proof}

\begin{theorem}[Three-vertex block norm bound]\label{thm:normblock}
Let $\|X\|_{2p}=\|Y\|_{2p}=1$, $1\le p<\infty$, and let
$\alpha,\beta,\eta,\delta,K\ge0$ satisfy
$\alpha\delta-\beta\eta\ge K^2$. No trace normalization is imposed.
For
\[
 T=\begin{pmatrix}
 \alpha XX^*+\beta YY^*&KXY^*\\
 KYX^*&\eta XX^*+\delta YY^*
 \end{pmatrix},\quad
 C=\begin{pmatrix}\alpha+\beta&K\\K&\eta+\delta\end{pmatrix},
\]
one has $\|T\|_p\le\|C\|_p$. The corresponding statement holds at
$p=\infty$ with operator-normalized $X,Y$. The constant one is sharp.
\end{theorem}
\begin{proof}
Write $U_0=\alpha+\beta$, $U_1=\eta+\delta$, and
$u=U_0U_1-K^2$. For positive row masses and $u>0$, use precisely the
vertices $A,B,C$ and weights $w_A,w_B,w_C$ in the proof of
Theorem~\ref{thm:block}. The condition is equivalent to
$\beta U_1+\eta U_0\le u$, so all three weights are nonnegative.
Their convex combination is the original coefficient tuple, while the
scalar comparator $C$ is unchanged at every vertex.

At vertex $A$, its factor is
\[
 M_A=\begin{pmatrix}(K/\sqrt{U_1})X&\sqrt{u/U_1}\,Y\\
                         \sqrt{U_1}\,Y&0\end{pmatrix}.
\]
The scalar matrix of its block $2p$-norms has Gram matrix $C$.
Permuting a block row or column reduces it to Lemma~\ref{lem:fangnorm}.
Since $\|M_AM_A^*\|_p=\|M_A\|_{2p}^2$, the desired bound follows.
At vertex $B$ use
\[
 M_B=\begin{pmatrix}\sqrt{U_0}\,X&0\\
                (K/\sqrt{U_0})Y&\sqrt{u/U_0}\,X\end{pmatrix};
\]
its compressed Gram matrix is again $C$. At the third vertex,
$T_C=\operatorname{diag}(X,Y)(C\otimes I_r)\operatorname{diag}(X,Y)^*$,
so Lemma~\ref{lem:altnorm} applies. The triangle inequality for the
Schatten norm combines these three bounds.

If $u=0$ and both row masses are positive, the coefficient condition
forces $\beta=\eta=0$ and the third-vertex argument applies directly.
If a row mass vanishes, $K=0$ and the remaining row is bounded by the
triangle inequality. These arguments include zero scalar coefficients.
The same proof uses the norm endpoints for $p=\infty$. Taking $X=Y$
to be a rank-one matrix of unit Schatten norm gives
$T=C\otimes XX^*$ and equality, proving sharpness.
\end{proof}

\begin{theorem}[Diagonal-filter multiplicativity]\label{thm:filterednorm}
Let $\Phi=\Phi_{a,b,\lambda}$ satisfy \eqref{eq:region}. For every
nonnegative diagonal $H$, every finite-dimensional completely positive
map $\Psi$, and every $1\le p\le\infty$,
\[
 \nu_p(\Phi_H\otimes\Psi)=\nu_p(\Phi_H)\nu_p(\Psi).
\]
The partner may have arbitrary input/output dimensions and need not
preserve trace.
\end{theorem}
\begin{proof}
For finite $p$, norm convexity reduces to a normalized pure input
$|0\rangle|u\rangle+|1\rangle|v\rangle$. Complete positivity yields a
factorization
\[
 (\operatorname{id}\otimes\Psi)(\psi)
 =\begin{pmatrix}XX^*&XY^*\\YX^*&YY^*\end{pmatrix}.
\]
Put $x=\|X\|_{2p}^2$, $y=\|Y\|_{2p}^2$, $m=x+y$ and $q=y/m$.
Suppose first $xy>0$, and normalize $X$ and $Y$ by $\sqrt{x}$ and
$\sqrt{y}$. The output before filtering, divided by $m$, has the
coefficients
\[
 \alpha=(1-q)(1-a),\quad\beta=qb,\quad
 \eta=(1-q)a,\quad\delta=q(1-b),\quad
 K=\lambda\sqrt{q(1-q)}.
\]
They obey
$\alpha\delta-\beta\eta-K^2=(k-\lambda^2)q(1-q)\ge0$.
For $H=\operatorname{diag}(h_0,h_1)$, multiply $\alpha,\beta$ by
$h_0^2$, $\eta,\delta$ by $h_1^2$, and $K$ by $h_0h_1$.
The determinant gap is multiplied by $h_0^2h_1^2$ and remains
nonnegative, including singular filters. Theorem~\ref{thm:normblock}
therefore gives
\[
 \|(\Phi_H\otimes\Psi)(\psi)\|_p
 \le m\|\Phi_H(\psi_q^+)\|_p,
 \quad |\psi_q^+\rangle=\sqrt{1-q}|0\rangle+\sqrt q|1\rangle.
\]
Finally,
\begin{align*}
 m&=\|\Psi(uu^*)\|_p+\|\Psi(vv^*)\|_p\\
 &\le(\|u\|^2+\|v\|^2)\nu_p(\Psi)=\nu_p(\Psi).
\end{align*}
If $m=0$ the output vanishes; if exactly one block norm is zero,
the input after $\operatorname{id}\otimes\Psi$ has a single basis
branch and the same estimate follows directly. Maximization proves the
upper bound and product inputs prove the reverse bound. The finite-dimensional
limit $p\to\infty$ gives the remaining endpoint.
\end{proof}

\subsection{Renyi radii and their scalar evaluation}
In the rest of this section logarithms are to base two. For $\alpha>1$
define
\[
 \widetilde\chi_\alpha(\mathcal N)=
 \inf_\sigma\sup_\rho\frac1{\alpha-1}
 \log_2\Tr\left[\left(\sigma^{\frac{1-\alpha}{2\alpha}}
 \mathcal N(\rho)\sigma^{\frac{1-\alpha}{2\alpha}}\right)^\alpha\right].
\]
The standard support convention gives infinity if an output is not
supported in $\sigma$. The ensemble-radius identity and product-ensemble
superadditivity used below are external results \cite{mo,beigi}.

\begin{corollary}\label{cor:renyiadd}
For every channel $\mathcal N$ and every $\alpha>1$,
\[
 \widetilde\chi_\alpha(\Phi\otimes\mathcal N)
 =\widetilde\chi_\alpha(\Phi)+\widetilde\chi_\alpha(\mathcal N).
\]
In particular, the $n$-use quantity is exactly
$n\widetilde\chi_\alpha(\Phi)$ for every positive integer $n$.
\end{corollary}
\begin{proof}
Phase covariance and data processing under discarding a uniform binary
phase flag show that dephasing a reference state cannot increase its
worst-output radius; this is the argument in \cite[Eq.~(37)]{fang}.
Thus diagonal references suffice for $\Phi$. Strictly positive references
suffice by approximation with $(1-\epsilon)\sigma+\epsilon I/d$;
the uniform radius penalty is at most $-\log_2(1-\epsilon)$.
For a fixed full-rank reference, the supremum is
$\alpha/(\alpha-1)$ times the logarithm of the maximum output
$\alpha$-norm of the corresponding sandwiched map.
Use a product reference for $\Phi\otimes\mathcal N$ and apply
Theorem~\ref{thm:filterednorm} with the sandwiched partner, which is a
completely positive map. Optimizing the two references proves the upper
bound. Product-ensemble superadditivity \cite{beigi}, with the
ensemble-radius identity \cite{mo}, proves the reverse bound.
\end{proof}

For a direct scalar formula set $\sigma_s=\operatorname{diag}(1-s,s)$,
$0<s<1$, $h_0=(1-s)^{(1-\alpha)/(2\alpha)}$,
$h_1=s^{(1-\alpha)/(2\alpha)}$, and $t=a+kq$. Define
\[
 B_{\alpha,s}(q)=\begin{pmatrix}
 h_0^2(1-t)&h_0h_1\lambda\sqrt{q(1-q)}\\
 h_0h_1\lambda\sqrt{q(1-q)}&h_1^2t
 \end{pmatrix}.
\]
If $L_\pm(\alpha,s,q)$ are its two eigenvalues, then
\begin{equation}\label{eq:renyiscalar}
 \widetilde\chi_\alpha(\Phi)=
 \inf_{0<s<1}\max_{0\le q\le1}
 \frac{\log_2(L_+^\alpha+L_-^\alpha)}{\alpha-1}.
\end{equation}
Indeed pure inputs suffice by norm convexity, and phase covariance
makes their eigenvalues depend only on $q$. No minimax interchange
or single-orbit optimal-ensemble assertion is made in this formula.

\subsection{The exact exponent above capacity}
Let $P_{\rm s}^{\star}(n,R;\Phi)$ be the best average success probability
for unassisted classical communication using $n$ copies, at least
$\lceil2^{nR}\rceil$ messages, arbitrary entangled input codewords and a
collective output measurement. The channel is memoryless; feedback and
preshared entanglement are excluded.

\begin{corollary}\label{cor:exponent}
For every $R\ge0$ the limit defining the optimal strong-converse exponent
exists and equals
\begin{equation}\label{eq:strongexponent}
 \lim_{n\to\infty}-\frac1n\log_2 P_{\rm s}^{\star}(n,R;\Phi)
 =\sup_{\alpha>1}\frac{\alpha-1}{\alpha}
       [R-\widetilde\chi_\alpha(\Phi)].
\end{equation}
The right side is positive for $R>C(\Phi)$. Every finite-block code obeys
the upper success bound obtained by exponentiating $-n$ times the right
side of \eqref{eq:strongexponent}. Product input codewords with collective
decoding attain the optimal asymptotic exponent.
\end{corollary}
\begin{proof}
The classical Arimoto converse bounds every code at each $\alpha>1$ by
$2^{-n(\alpha-1)(R-\widetilde\chi_\alpha(\Phi))/\alpha}$ after
Corollary~\ref{cor:renyiadd}. Optimizing $\alpha$ gives the finite-block
statement. Equality of the asymptotic exponent follows from the general
regularized exponent theorem \cite[Theorem~5 and Appendix~E]{ct}, since
Corollary~\ref{cor:renyiadd} removes its regularization at every order.
For product-codeword attainment, apply the classical-quantum coding
exponent theorem \cite[Section~VI]{mo} to the channel with alphabet all
pure qubit input states. Its optimized radius equals that of $\Phi$
by norm convexity; the finite-dimensional output permits finite
approximating alphabets. This gives the same exponent with product
codewords. This coding theorem is an external operational dependency,
not a construction of a particular finite-block decoder in this paper.
\end{proof}

The new norm proof extends Fang's GAD boundary to the full region
$0\le\lambda^2\le k$. The standard unital, entanglement-breaking and
zero-temperature overlaps are not claimed as new. For example,
$a=1/10,b=3/10,\lambda^2=2/5$ is nonunital, has strict extra dephasing
because $2/5<k=3/5$, and is not entanglement breaking because
$\lambda^2>ab$. The CP-allowed strip $k<\lambda^2\le k+ab$ is not
covered. Neither the general two-row norm-compression conjecture nor
any qutrit-ladder capacity statement follows from these results.
Integer-order filtered multiplicativity is already covered by the
entrywise-positive-map theorem of King, Nathanson and Ruskai
\cite{knr}. Without filtering, the general qubit results of King and
Koldan \cite{koldan} already cover $p=2$ and $p\ge4$.
The extension here must therefore be read with these overlaps, in
particular the real noninteger orders and the filtered nonunital,
non-entanglement-breaking interior, rather than as a new result for
every individual parameter and order in the theorem.
